\section{Introduction}

\subsection{Study Rationale}

Pancreatic ductal adenocarcinoma (PDAC) is the third leading cause of cancer
death in the United States and the fourth in the European Union, with total
five-year overall survival below 13 percent \cite{Siegel2025}. The
pancreaticoduodenectomy, or Whipple procedure, is the most technically demanding
of all curative-intent abdominal oncology operations: it requires resection and
reconstruction of four luminal structures (the pancreas, the duodenum, the common
bile duct, and the stomach or pylorus) and dissection around two named vessels
(the superior mesenteric vein and the portal vein), with an artery-first approach
to the superior mesenteric artery. The leading lethal postoperative sequence is
the fistula-sepsis cascade: pancreatic fistula, intra-abdominal infection,
hemorrhage, sepsis, and failure to rescue, graded under the ISGPS classification
of postoperative pancreatic fistula \cite{Bassi2017}.

The 2025 Dutch nationwide cohort of 1000 robotic pancreaticoduodenectomies
reported a conversion rate of 10.1 percent, Clavien-Dindo grade III or higher
complications of 41.3 percent, ISGPS grade B or C fistula of 24.4 percent, and
30-day mortality of 3.9 percent \cite{DutchCohort2025}; these figures anchor the
contemporary human baseline. The Phase 1 predicate \cite{kawchak2026phase1}
established the recommended Phase 2 dose (RP2D, 300 mg daraxonrasib once daily)
and the feasibility and safety of the on-premises large language model (LLM)
directed eight-arm robotic Whipple, but a single-arm feasibility study cannot
establish efficacy. That predicate result creates genuine clinical equipoise
between the integrated combination and the prevailing standard of care, and it is
precisely that equipoise that this randomized, controlled efficacy comparison is
designed to resolve. The present Phase 2 protocol therefore randomizes patients
1:1 (open-label, with blinded independent central review, BICR) between Arm A,
perioperative daraxonrasib at RP2D plus the on-premises LLM-directed robotic
Whipple (PancreSpeed II), and Arm B, modified FOLFIRINOX \cite{Conroy2018} plus a
standard high-volume pancreaticoduodenectomy.

The clinical argument for the combination, and for funding a definitive test of
it, is best stated through four counterfactual scenarios, summarized in Figure 3,
in which withholding the integrated approach shortens progression-free survival
(PFS) and overall survival (OS) or, in the fourth case, denies the field a
trustworthy answer altogether. In Scenario A, a borderline-resectable tumor
abutting the superior mesenteric vein progresses to unresectable disease during a
human-only scheduling delay (operating room access, fatigue, staging lag), and the
R0 window closes; the LLM-directed robot instead achieves a rapid, precise,
in-window R0 resection. In Scenario B, an unrecognized superior-mesenteric or
portal-vein injury during manual surgery triggers the
hemorrhage-conversion-fistula-sepsis-failure-to-rescue cascade; the five-vessel
no-fly vascular gate with a $\leq 3$ ms cross-arm emergency stop averts the injury
and preserves a complete resection. 

In Scenario C, manual mistiming of the
daraxonrasib restart (too early or too late) causes anastomotic dehiscence or
micrometastatic outgrowth; the LLM advisory restart at T+7, T+14, or T+21 days,
keyed to fistula status and drug trough, preserves the therapeutic window. In
Scenario D, an under-funded, single-center, under-powered study fails to deliver a
definitive, generalizable, and equitable answer at all: slow accrual, attrition,
and a thin biomarker core leave the central PFS question unresolved for the next
patient. The patient-aligned co-investment model converts capital, strictly behind
a structural firewall, into accrual speed, retention, fidelity, and statistical
power, raising the probability of a definitive result without ever touching the
endpoints, their adjudication, or the analysis.

These scenarios invite four framing questions that this protocol is designed to
answer constructively rather than rhetorically. First, given the narrow resection
window and the unforgiving fistula-sepsis cascade, how did we ever rely on humans
alone for this treatment, with no second-opinion oracle isolated from the
operative kinematics. Second, why did we fully trust a prior trial system with so
many opportunities for compounding human error across scheduling, dissection, and
drug timing, when each step admits a documented, deterministic safeguard. Third,
as AI outputs grow more complex and voluminous, why would human peer review alone
continue to suffice for increasingly intricate AI-directed care, rather than a
verification-before-generation discipline with a hash-chained audit trail
\cite{kawchak2026hr9510}. Fourth, when aligned philanthropic and private capital
can raise the success likelihood of a definitive trial without altering a single
endpoint, why would we leave that answer under-powered and inequitable. The
protocol answers the first three by retaining the human in the loop while making
every machine action validated, bounded, and re-traceable, and it answers the
fourth with a Patient-Aligned Co-Investment Facility whose capital firewall lets
contributions buy only operational success-likelihood, never a favorable result
\cite{kawchak2026hr9510}.

\begin{mermaidfig}
% Scenario A
\node[mmin]  (a0) at (0,0)       {Scenario A\\Borderline PDAC,\\SMV abutment $75^\circ$};
\node[mmdark](aw) at (5.2,0.95)  {Human-only delay:\\unresectable, R0\\window closes, shorter PFS/OS};
\node[mmgoal](ac) at (5.2,-0.95) {LLM plus robot:\\in-window R0 resection};
% Scenario B
\node[mmin]  (b0) at (0,-3.7)    {Scenario B\\Intra-op SMV/PV plane};
\node[mmdark](bw) at (5.2,-2.75) {Manual injury: hemorrhage,\\conversion, R1/R2,\\fistula-sepsis cascade};
\node[mmgoal](bc) at (5.2,-4.65) {No-fly gate, $\leq 3$ ms\\e-stop: injury averted};
% Scenario C
\node[mmin]  (c0) at (0,-7.4)    {Scenario C\\Perioperative\\daraxonrasib};
\node[mmdark](cw) at (5.2,-6.45) {Manual mistiming:\\dehiscence or\\micrometastatic outgrowth};
\node[mmgoal](cc) at (5.2,-8.35) {Advisory T+7/14/21:\\optimally timed restart};
% Scenario D
\node[mmin]  (d0) at (0,-11.1)   {Scenario D\\Funding the\\definitive answer};
\node[mmdark](dw) at (5.2,-10.15){Under-funded, single-center,\\under-powered};
\node[mmgoal](dc) at (5.2,-12.05){Co-invested behind firewall: power, retention, equity};

\draw[mmedge] (a0.east) -- node[above, font=\scriptsize, align=center, xshift=-0.075cm, yshift=0.315cm]{human only} (aw.west);
\draw[mmedge] (a0.east) -- node[below,font=\scriptsize, align=center, xshift=-0.075cm, yshift=-0.375cm]{combination} (ac.west);

\draw[mmedge] (b0.east) -- node[above,font=\scriptsize, align=center, xshift=-0.075cm, yshift=0.315cm]{human only} (bw.west);
\draw[mmedge] (b0.east) -- node[below,font=\scriptsize, align=center, xshift=-0.075cm, yshift=-0.375cm]{combination} (bc.west);

\draw[mmedge] (c0.east) -- node[above,font=\scriptsize, align=center, xshift=-0.075cm, yshift=0.315cm]{human only} (cw.west);
\draw[mmedge] (c0.east) -- node[below,font=\scriptsize, align=center, xshift=-0.075cm, yshift=-0.375cm]{combination} (cc.west);

\draw[mmedge] (d0.east) -- node[above,font=\scriptsize, align=center, xshift=-0.075cm, yshift=0.425cm]{under-funded} (dw.west);
\draw[mmedge] (d0.east) -- node[below,font=\scriptsize, align=center, xshift=-0.075cm, yshift=-0.375cm]{co-invested} (dc.west);
\end{mermaidfig}
\vspace{-0.7cm}
\figcaption{Figure 3. The four counterfactual scenarios in which withholding the
integrated approach worsens the patient or\\the evidence. Scenario A: a human-only
scheduling delay lets a borderline-resectable tumor progress and closes the\\R0
window. Scenario B: an unrecognized superior-mesenteric or portal-vein injury
during manual surgery triggers\\the fistula-sepsis cascade; the no-fly gate averts it. Scenario C: manual mistiming of the
daraxonrasib restart\\causes dehiscence or micrometastatic outgrowth; the timed
advisory prevents it. Scenario D: an under-funded, single-center, under-powered
study yields no definitive equitable answer; patient-aligned co-investment, walled
off from endpoints by capital firewall, converts capital into accrual
speed, retention, fidelity, and statistical power.}

\subsection{Background}

\subsubsection*{PDAC epidemiology and the Whipple fistula-sepsis cascade}

The epidemiologic burden and operative hazard of PDAC define the unmet need this
protocol addresses. With under-13-percent five-year survival \cite{Siegel2025}
and a most-demanding curative operation whose leading complication is a
self-amplifying fistula-sepsis-hemorrhage cascade \cite{Bassi2017}, the
population has few alternatives and a narrow margin for technical error. The
contemporary robotic baseline \cite{DutchCohort2025} demonstrates that even
expert, high-volume robotic surgery leaves a substantial residual of conversions,
high-grade complications, and clinically relevant fistulae. The combination in
this protocol targets precisely the failure modes that the baseline still
tolerates, and the randomized internal control (Arm B) measures the increment
against modified FOLFIRINOX plus standard pancreaticoduodenectomy
\cite{Conroy2018}.

\subsubsection*{Daraxonrasib mechanism and clinical program}

Daraxonrasib (RMC-6236) is an oral RAS(ON) multi-selective inhibitor that binds
the active, GTP-bound (ON) state of multiple RAS variants, providing broad
pan-KRAS coverage including the G12 alleles that predominate in PDAC. The FDA
granted daraxonrasib Breakthrough Therapy designation in June 2025 for previously
treated metastatic PDAC with KRAS G12 mutations \cite{revbreakthrough}, and the
late-phase program comprises RASolute 302 (a Phase 3 study in previously treated
metastatic PDAC) and RASolve 301 (a first-line program) \cite{rasolute302}. The
broad pan-KRAS eligibility of RASolute 302 defines the molecular eligibility used
in this protocol, and the perioperative pause-and-restart logic respects any
RASolve 301 checkpoint-inhibitor combination context. Five prior computational
works by the author independently complemented this drug-development program
throughout 2025 and 2026 and supply the simulation scaffolding adapted here: a
60-second PDAC robotic Whipple and daraxonrasib simulation
\cite{kawchak_2026_20196639}; an AI digital-twin pancreatic-cancer simulation for
in silico FDA-compliance and cost efficiency \cite{kawchak_2025_17239510}; a
quantitative systems pharmacology (QSP) metastatic-PDAC trial simulation with code,
VVUQ, and playbook \cite{kawchak_2025_17001137}; a 100,000-patient in silico
Phase 3 five-arm PDAC trial triplicate \cite{kawchak_2025_16415815}; and
end-to-end PDAC digital-twin clinical-trial proposals \cite{kawchak_2025_15735068}.

\subsubsection*{Current state of AI in clinical trials and surgery}

Artificial intelligence in regulated medicine has, to date, been narrow and
prediction-oriented rather than generative and intervention-directing. As of the
current taxonomy of FDA-authorized medical devices, AI is overwhelmingly deployed
as prediction-only or narrow task-specific software, with reported authorizations
spanning 1,016 in the published taxonomy \cite{Singh2025} and rising to 1,450 or
more by current counts, while zero generative LLM devices have been approved for
patient-facing intervention. This protocol therefore does not propose a generative
device; it proposes a narrow, locked, validated device behavior governed by an
on-premises repository-based LLM whose advisories are bounded, auditable, and
held outside the robot vendor kinematic stack. The author's established record of
LLM evaluation and trust over the period from August 2024 to June 2026 grounds
this design, including a comparative code-generation evaluation of 16 proprietary
versus open-source LLMs against the FDA Adverse Event Reporting System
\cite{kawchak2025codegen16}, a mobile pancreatic-cancer surgical-humanoid program
with priority VVUQ \cite{kawchak2026unitree}, the 21 CFR part 312 Physical AI
unification adaptation \cite{kawchak2026adapt312}, and the
verification-before-generation legislative framing \cite{kawchak2026hr9510}.

\subsubsection*{Funding, co-investment, and the success likelihood of a definitive
trial}

A definitive, generalizable, and equitable Phase 2 answer is itself an outcome
that can be made more or less likely by how the trial is resourced, and this
protocol makes that mechanism explicit and governed. The Patient-Aligned
Co-Investment Facility (PACIF) allows wealthier individuals connected to cancer
patients to contribute ring-fenced philanthropic or private capital with the sole
purpose of raising the trial's success likelihood. The defining safeguard is a
capital firewall: contributors are structurally walled off from randomization,
endpoint definition, outcome adjudication, statistical analysis, and publication,
so that capital can never buy a favorable result.

What capital can buy is limited
to six operational success-likelihood levers, each of which acts on accrual,
fidelity, retention, data quality, or oversight rather than on the answer. First,
multi-site expansion to eight academic hepato-pancreato-biliary centers
accelerates accrual and broadens generalizability. Second, Phase 0 compute funds
$\geq 5000$ simulations across $\geq 3$ independent frameworks, lowering the
sim-to-real gap (target $< 1.5$ mm and $< 0.4$ N) and raising the unified safety
level to USL $\geq 8.0$. Third, a robot fleet with redundancy increases throughput
and uptime. Fourth, a Patient Access and Equity Fund covers travel, lodging,
lost-wage, caregiver, and navigation costs, raising retention, which lowers
dropout, and protects statistical power, while making enrollment more
equitable and representative. 

Fifth, central infrastructure (BICR, central
pathology, ctDNA, federated learning, and biostatistics) improves data quality and
reduces bias. Sixth, independent oversight (a Data and Safety Monitoring Board, an
independent safety monitor, a Physical AI Safety Review Committee, and an
ethicist) protects safety and integrity. The facility is governed by the H.R. 9510
verification, validation, and uncertainty quantification (VVUQ) financial-data
standard \cite{kawchak2026hr9510}, disclosed under 21 CFR part 54
\cite{ecfr2026part54}, and its results are deposited openly under CC BY 4.0.
Table~\ref{tab:coinvest} maps each co-investment tranche to the operational lever
it funds and to the firewall that keeps it away from the answer. In short, aligned
capital raises power, fidelity, retention, equity, and generalizability without
buying the answer.

\subsubsection*{Physical AI concerns addressed by this protocol}

Nine recurring concerns are raised about Physical AI in patient care, and this
protocol pairs each with a concrete, prespecified mitigation. The first eight
concern the device and its governance; the ninth, new to this Phase 2 design,
concerns financial influence and inequitable access and is answered by the capital
firewall and the equity fund. Figure 4 maps the nine concerns to their
mitigations, and Table~\ref{tab:concerns} states each pairing in full.


\subsection{Risk/Benefit Assessment}

\subsubsection{Known Potential Risks}

The known potential risks fall into three categories. Drug risks arise from
daraxonrasib toxicity, including gastrointestinal effects (nausea, diarrhea,
stomatitis), dermatologic and mucocutaneous events, hepatic enzyme elevations,
and dose-limiting toxicities bounded at the established RP2D; perioperative
exposure additionally raises a theoretical risk to wound and anastomotic healing
if restart timing is suboptimal. Device and robotic risks include device
malfunction, an erroneous AI advisory, cyber-physical failure of the control or
telemetry chain, and inadvertent vascular injury; these are the categories that
motivate the significant-risk determination and the layered hard safety limits.
Operative risks are those intrinsic to the Whipple procedure: postoperative
pancreatic fistula (graded by ISGPS \cite{Bassi2017}), hemorrhage, delayed gastric
emptying, intra-abdominal infection and sepsis, and the need for conversion to
open surgery. Because the design is randomized, patients in the control arm (Arm B)
also carry the operative and modified-FOLFIRINOX toxicity risks intrinsic to the
standard of care \cite{Conroy2018}. The contemporary robotic baseline quantifies residual magnitude of these operative risks even in expert hands
\cite{DutchCohort2025}.



\begin{mermaidfig}
\node[mmin]  (c1) at (0,0)       {Physical AI\\limitations};
\node[mmin]  (c2) at (0,-1.6)    {Patient safety\\is paramount};
\node[mmin]  (c3) at (0,-3.2)    {Loss of\\human workers};
\node[mmin]  (c4) at (0,-4.8)    {Single-source\\software};
\node[mmin]  (c5) at (0,-6.4)    {Proprietary-model\\dependency};
\node[mmin]  (c6) at (0,-8.0)    {Open-source LLMs\\not domestic};
\node[mmin]  (c7) at (0,-9.6)    {Overly complex\\AI workflows};
\node[mmin]  (c8) at (0,-11.2)   {LLMs are\\black boxes};
\node[mmin]  (c9) at (0,-12.8)   {Financial influence,\\inequitable access};
\node[mmgoal](m1) at (4.6,0)      {Phase 0 sim validation;\\narrow locked device};
\node[mmgoal](m2) at (4.6,-1.6)   {10 kHz heartbeat, $\leq 3$ ms\\e-stop, force caps};
\node[mmgoal](m3) at (4.6,-3.2)   {Human-in-the-loop; surgeon\\plus backup plus ISM};
\node[mmgoal](m4) at (4.6,-4.8)   {Cross-framework\\($\geq 3$ frameworks)};
\node[mmgoal](m5) at (4.6,-6.4)   {On-premises LLM,\\open weights, isolated};
\node[mmgoal](m6) at (4.6,-8.0)   {Domestic deployment;\\data residency preserved};
\node[mmgoal](m7) at (4.6,-9.6)   {Deny-by-default,\\simplified MCP gate};
\node[mmgoal](m8) at (4.6,-11.2)  {Hash-chained audit\\trail; 21 CFR part 11};
\node[mmgoal](m9) at (4.6,-12.8)  {Capital firewall, equity fund,\\H.R. 9510};
\draw[mmedge] (c1) -- (m1);
\draw[mmedge] (c2) -- (m2);
\draw[mmedge] (c3) -- (m3);
\draw[mmedge] (c4) -- (m4);
\draw[mmedge] (c5) -- (m5);
\draw[mmedge] (c6) -- (m6);
\draw[mmedge] (c7) -- (m7);
\draw[mmedge] (c8) -- (m8);
\draw[mmedge] (c9) -- (m9);
\end{mermaidfig}
\vspace{-0.65cm}
\figcaption{Figure 4. The mapping of all nine Physical AI concerns to protocol
mitigations, the same nine pairs in Table~\ref{tab:concerns}.\\ The original eight
device concerns are paired with validated device behavior over open-endedness,
layered hard\\safety limits, retained human roles, and multi-framework
non-single-source validation; the ninth, financial\\influence and inequitable
access, is answered by the capital firewall, the Patient Access and Equity Fund,\\21 CFR part 54 disclosure, the H.R. 9510 VVUQ financial standard, and open CC BY
deposition.}

\vspace{0.3em}

\begin{table}[t]\centering\footnotesize
\caption{The nine Physical AI concerns and their protocol mitigations.}
\label{tab:concerns}
\begin{tabularx}{\textwidth}{L{5.5cm} Y}
\toprule
\textbf{Concern} & \textbf{Protocol mitigation} \\
\midrule
Physical AI limitations & Phase 0 simulation validation over $\geq 5000$ procedures and $\geq 3$ independent frameworks; device behavior is narrow and locked, not generative, and operates only within validated envelopes. \\
Patient safety is paramount & Layered hard limits: a 10 kHz heartbeat broadcast bus, a $\leq 3$ ms cross-arm emergency stop, $\leq 3$ N per-arm and $\leq 18$ N cumulative tip-force caps, and a five-vessel no-fly gate around the SMV and PV. \\
Loss of human workers & The human is retained in the loop: a credentialed surgeon, a backup operator, and an independent safety monitor (ISM) supervise every procedure, with the LLM acting as a bounded advisory oracle rather than a replacement. \\
Single-source software & Validation is cross-framework, with $\geq 3$ independent frameworks required, so no single simulator or codebase is a single point of failure. \\
Proprietary-model dependency & The advisory uses an on-premises repository-based LLM with open weights, isolated from any single vendor stack, so the protocol is not captive to one proprietary model provider \cite{kawchak2025codegen16}. \\
Non-domestic open-source LLMs & Deployment is domestic and on-premises, preserving data residency and avoiding reliance on non-domestic open-source model hosting; inference and records remain within the controlled environment. \\
Overly complex AI workflows & A deny-by-default posture with a deliberately simplified gate surface (a small set of Model Context Protocol servers) constrains the action space and reduces workflow complexity to an auditable minimum. \\
LLMs are black boxes & A hash-chained audit trail under 21 CFR part 11 records every advisory, command, and safety event, making each re-traceable to a deterministic seed and a specific commit, so the black box becomes traceable \cite{kawchak2026adapt312}. \\
Financial influence and inequitable access & A capital firewall walls funders off from randomization, endpoints, adjudication, analysis, and publication; a Patient Access and Equity Fund supports equitable enrollment and retention; influence is disclosed under 21 CFR part 54, governed by the H.R. 9510 VVUQ financial standard, and all results are deposited openly under CC BY \cite{kawchak2026hr9510,ecfr2026part54}. \\
\bottomrule
\end{tabularx}
\end{table}

\begin{table}[t]\centering\footnotesize
\caption{The patient-aligned co-investment tranches, the operational
success-likelihood levers they fund, and the capital firewall.}
\label{tab:coinvest}
\begin{tabularx}{\textwidth}{L{3.2cm} L{4.3cm} Y}
\toprule
\textbf{Co-investment tranche} & \textbf{Operational lever funded} & \textbf{Effect on success likelihood; firewall note} \\
\midrule
Multi-site network (8 centers) & Eight academic HPB centers activated and supported & Faster accrual and broader generalizability of the PFS result; no influence on randomization or endpoints. \\
Phase 0 compute & $\geq 5000$ simulations across $\geq 3$ frameworks & Lower sim-to-real gap ($< 1.5$ mm, $< 0.4$ N); USL $\geq 8.0$, raising fidelity; compute cannot alter the analysis plan. \\
Robot fleet plus redundancy & Device fleet, spares, and uptime engineering & Higher throughput and uptime, sustaining in-window accrual; no access to outcome adjudication. \\
Patient Access and Equity Fund & Travel, lodging, lost-wage, caregiver, and navigation support & Higher retention that lowers dropout and protects statistical power, with more equitable and representative enrollment; funders see no patient-level outcomes. \\
Central review and biomarker core & BICR, central pathology, ctDNA, and federated learning & Improved data quality and reduced bias; the core is independent of contributors and of the analysis team. \\
Independent oversight & DSMB, ISM, Physical AI Safety Review Committee, ethicist & Stronger safety and integrity monitoring; oversight bodies are insulated from funders. \\
\bottomrule
\end{tabularx}
\end{table}





\vspace{-0.2cm}

\subsubsection{Known Potential Benefits}

The known potential benefits are directed at the specific failure modes above. A
precise, in-window R0 resection is the principal oncologic benefit, addressing the
resection-window collapse and the margin-positivity that drive shorter survival.
Layered hard safety limits with a full hash-chained audit trail provide a benefit
not available in manual surgery: every machine action is bounded and re-traceable,
supporting both intraoperative safety and post hoc verification. Optimized
drug-restart timing, advised at T+7, T+14, or T+21 days and keyed to fistula
status and drug trough, aims to restore systemic therapy at the earliest safe
point and thereby suppress micrometastatic outgrowth. Beyond the individual
procedure, the Patient Access and Equity Fund confers a direct participant
benefit, lowering the travel, lodging, lost-wage, and caregiver burdens that drive
attrition and inequitable enrollment, which improves both retention and the
representativeness of the evidence. For a low-survival population with few
alternatives, even incremental gains in margin status, complication avoidance,
therapy continuity, and equitable access are clinically meaningful.


\textit{Capital firewall}: contributors to any tranche cannot influence randomization, endpoint definition, outcome adjudication, statistical analysis, or publication; capital buys only operational success-likelihood, never a favorable result. Influence is disclosed under 21 CFR part 54 and governed by the H.R. 9510 VVUQ financial-data standard \cite{kawchak2026hr9510,ecfr2026part54}.

\vspace{-0.2cm}

\subsubsection{Assessment of Potential Risks and Benefits}

Figure 5 frames the benefit-risk weighing for this population, now as a randomized
comparison with an internal control. Against a baseline of under-13-percent
five-year survival and few alternatives, the novel risks (device malfunction,
advisory error, cyber-physical failure) are weighed against the benefits (precise
in-window R0 resection, layered hard safety limits with a full audit trail,
optimized drug-restart timing, and the access-fund gains in retention and equity).
For this low-survival, limited-alternative population, the benefit-risk balance is
favorable: the potential benefits are likely to outweigh the risks, and the
assessment is consistent with the expanded-access ethic that permits
investigational treatment for serious or immediately life-threatening conditions
when the potential benefit justifies the potential risk (21 CFR \S312.84).

\begin{mermaidfig}
\node[mmdark](base) at (0,-0.95)  {Advanced PDAC baseline\\5-year OS $<$ 13\%\\3rd US / 4th EU death};
\node[mmin]  (r1)  at (-5.1,-2.4) {Device malfunction};
\node[mmin]  (r2)  at (-5.1,-4.4) {AI advisory error};
\node[mmin]  (r3)  at (-5.1,-6.4) {Cyber-physical failure};
\node[mmstep](b1)  at (5.1,-2.4)  {In-window R0 resection};
\node[mmstep](b2)  at (5.1,-4.4)  {Hard limits plus audit trail};
\node[mmstep](b3)  at (5.1,-6.4)  {Restart timing plus equity fund};
\node[mmdec] (w)   at (0,-4.4)    {Benefit-risk for\\low-survival patients};
\node[mmgoal](j)   at (0,-7.6)    {Favorable: benefits likely outweigh\\risks; \S312.84};
\draw[mmedge] (base) -- (w);
\draw[mmedge] (r1.east) -- (w.west);
\draw[mmedge] (r2) -- (w);
\draw[mmedge] (r3.east) -- (w.west);
\draw[mmedge] (b1.west) -- (w.east);
\draw[mmedge] (b2) -- (w);
\draw[mmedge] (b3.west) -- (w.east);
\draw[mmedge] (w) -- (j);
\end{mermaidfig}
\vspace{-0.45cm}
\figcaption{Figure 5. Risk-benefit framework for the randomized Phase 2
comparison in advanced PDAC. Against a baseline\\of under-13-percent five-year
survival and few alternatives, the novel risks (device malfunction, advisory
error, cyber-physical failure) are weighed against the benefits (precise in-window
R0 resection, layered hard safety\\limits with a full audit trail, optimized drug
timing, and the access-fund gains in retention and equity).\\For this low-survival,
limited-alternative population the benefit-risk balance is favorable.}

Finally, Figure 6 makes the co-investment logic explicit and shows where the
firewall sits. Aligned capital and the PACIF supply only the operational levers;
the capital firewall blocks every path from a contributor to the endpoints, their
adjudication, or the analysis, so the levers raise the success likelihood of a
definitive answer while the integrity of that answer remains structurally
protected.

\begin{mermaidfig}
\node[mmin]  (cap) at (0,0.5)       {Aligned capital\\and PACIF tranches};
\node[mmdec] (fw)  at (0,2.55)     {Capital\\firewall};
\node[mmstep](l1)  at (6,2.55)  {8 academic\\HPB sites};
\node[mmstep](l2)  at (6,0.85)  {Phase 0 compute\\($\geq 5000$ sims)};
\node[mmstep](l3)  at (6,-0.85) {Patient Access\\and Equity Fund};
\node[mmstep](l4)  at (6,-2.55) {Central review\\core (BICR)};
\node[mmgoal](o1)  at (10.2,2.55) {Higher\\statistical power};
\node[mmgoal](o2)  at (10.2,0.85) {Lower\\sim-to-real gap};
\node[mmgoal](o3)  at (10.2,-0.85){Lower dropout,\\equitable enrollment};
\node[mmgoal](o4)  at (10.2,-2.55){Definitive, generalizable,\\equitable answer};
\node[mmdark](block) at (0,-2.55){Endpoints, adjudication,\\analysis: no funder};
\draw[mmedge] (cap) -- (fw);
\draw[mmedge] (fw.east) -- (l1.west);
\draw[mmedge] (fw.east) -- (l2.north west);
\draw[mmedge] (fw.east) -- (l3.north west);
\draw[mmedge] (fw.east) -- (l4.north west);
\draw[mmedge] (l1) -- (o1);
\draw[mmedge] (l2) -- (o2);
\draw[mmedge] (l3) -- (o3);
\draw[mmedge] (l4) -- (o4);
\draw[mmedged, rounded corners=4pt] (fw.west) -- ++(-0.75,0) |- node[pos=0.25, right, font=\scriptsize, align=center, yshift=-1.15cm]{blocked} (block.west);
\end{mermaidfig}
\vspace{-0.7cm}
\figcaption{Figure 6. The co-investment-to-success-likelihood pathway. Aligned
capital and the PACIF flow only through the capital firewall into operational
levers (eight sites, Phase 0 compute, the Patient Access and Equity Fund, and the
central review core), which raise statistical power, lower the sim-to-real gap,
lower dropout, broaden equitable enrollment, and so raise the probability of a
definitive, generalizable, equitable answer. The dashed blocked edge shows that
the firewall admits no path from funders to the endpoints, their adjudication, or
the analysis: capital raises operational success likelihood while the integrity of
the result is structurally protected \cite{kawchak2026hr9510,ecfr2026part54}.}
